\section{Geometric phase, Berry/Wilczek--Zee connections, and gauge structure}
\label{sec:geophase}

\subsection{Ordered narrowband waves and polarization transport}

We assume the substrate admits an ordered, narrowband carrier sector that supports multi-component polarized
waves with a long phase-coherence length.
Locally, a narrowband excitation can be written as
\begin{equation}
  \vecu(x,t)\;\approx\;\mathrm{Re}\Big\{ a(x,t)\,\bm{\epsilon}(x,t)\,e^{\ii\theta(x,t)} \Big\},
  \qquad
  \theta(x,t)=\phi(x,t)-\omega_0 t,
  \label{eq:quasi-monochromatic}
\end{equation}
where $a(x,t)\ge 0$ varies slowly, $\bm{\epsilon}(x,t)\in\C^{N}$ is a normalized polarization state
($\bm{\epsilon}^\dagger\bm{\epsilon}=1$), and $\omega_0$ is the carrier angular frequency.
The internal dimension $N$ counts the relevant components (e.g.\ embedding-direction components and/or a
coarse-grained mode basis).

A central point is that polarization transport is not unique: the same physical state can be represented in many
local phase conventions.
Geometric phase formalizes this redundancy as a connection on a bundle of polarization states.

\subsection{Rank-1 Berry connection (U(1))}

For a single isolated polarization mode, represent the local state by a normalized vector
$\lvert u(x)\rangle\in\C^{N}$ that depends smoothly on spacetime position $x^\mu$
(here $x^\mu=(c\,t,x^1,x^2,x^3)$ built from the intrinsic coordinates and a chosen characteristic speed $c$ of the
carrier branch).
Define the Berry connection
\begin{equation}
  a_\mu(x) := \ii\,\langle u(x)\,|\,\partial_\mu u(x)\rangle.
  \label{eq:berry-connection}
\end{equation}
Under a local phase choice $\lvert u(x)\rangle\mapsto e^{-\ii\chi(x)}\lvert u(x)\rangle$ one finds
\begin{equation}
  a_\mu \mapsto a_\mu + \partial_\mu \chi,
  \label{eq:berry-gauge}
\end{equation}
so only gauge-invariant quantities carry physical meaning.
The Berry curvature is the field strength
\begin{equation}
  f_{\mu\nu} := \partial_\mu a_\nu - \partial_\nu a_\mu,
  \label{eq:berry-curvature}
\end{equation}
and the holonomy (geometric phase) around a closed loop $\Gamma$ is
\begin{equation}
  \gamma_\Gamma
  = \oint_\Gamma a_\mu\,\dd x^\mu
  = \int_{\Sigma:\,\partial\Sigma=\Gamma} f_{\mu\nu}\,\dd\Sigma^{\mu\nu}.
  \label{eq:berry-holonomy}
\end{equation}

\subsection{What the geometric phase is ``of'': eigenmode bundles and adiabaticity}
\label{subsec:berry-of}

The connection above is not attached to an arbitrary decomposition of a real field into amplitude and phase.
In the present setting one should choose $\lvert u(x)\rangle$ as a (locally defined) eigenvector of the
linearized carrier operator about the slowly varying background configuration at $x$.
As $x$ varies, these eigenvectors form an eigenbundle, and the Berry/Wilczek--Zee connection is the natural
connection on that bundle.

A clean rank-1 ($U(1)$) reduction requires that the dynamics remains in a single isolated band.
If $\Delta\omega_{\mathrm{gap}}(x)$ is the local separation to the nearest competing band and
$\Omega_{\mathrm{mix}}(x)$ is the rate of non-adiabatic mode conversion induced by background gradients,
then the adiabatic condition
\begin{equation}
  \Omega_{\mathrm{mix}}(x)\ll \Delta\omega_{\mathrm{gap}}(x)
  \label{eq:adiabatic-condition}
\end{equation}
should hold on the envelope scale.
When this fails, a measured ``Berry signal'' typically reflects mode mixing or beating, and the correct reduction
is to retain the relevant multi-dimensional subspace (the Wilczek--Zee description).
In the brane model, the pre-stress parameter $\alpha$ and the wavelength regime control branch separation and
therefore how well an Abelian gauge description can apply.

\subsection{Wilczek--Zee connection (non-Abelian)}

If the relevant polarization sector is an $n$-dimensional near-degenerate subspace, choose an orthonormal
frame $\mathbf{E}(x)\in\C^{N\times n}$ with columns $\lvert u_a(x)\rangle$ spanning the subspace, so that
$\mathbf{E}^\dagger\mathbf{E}=I_n$.
Define the Wilczek--Zee connection
\begin{equation}
  \mathcal{A}_\mu(x) := \ii\,\mathbf{E}(x)^\dagger\,\partial_\mu \mathbf{E}(x)\in\mathfrak{u}(n).
  \label{eq:wz-connection}
\end{equation}
A change of local frame $\mathbf{E}\mapsto \mathbf{E}\,U(x)$ with $U(x)\in U(n)$ yields the non-Abelian gauge law
\begin{equation}
  \mathcal{A}_\mu \mapsto U^\dagger \mathcal{A}_\mu U + \ii\,U^\dagger \partial_\mu U,
  \label{eq:wz-gauge}
\end{equation}
and the curvature (field strength) is
\begin{equation}
  \mathcal{F}_{\mu\nu}
  := \partial_\mu \mathcal{A}_\nu-\partial_\nu \mathcal{A}_\mu
     -\ii\,[\mathcal{A}_\mu,\mathcal{A}_\nu].
  \label{eq:wz-curvature}
\end{equation}

\subsection{Axis-aligned triplet phases and coupling-induced mixing ($U(1)^3\rightarrow U(3)$)}
\label{subsec:axis-triplet}

On a cubic lattice substrate it is natural to posit a three-component narrowband carrier triplet
$\Psi=(\psi_x,\psi_y,\psi_z)^\top\in\C^3$ aligned with the lattice axes.
Introduce a dimensionless mixing parameter $g_{\mathrm{mix}}\ge 0$ controlling inter-axis mixing.
In the lattice substrate, $g_{\mathrm{mix}}$ is determined primarily by the prestretch $\alpha$, because prestretch converts transverse amplitude into lateral restoring forces that couple the nominal axis channels; the microscopic coupling stencil (notably off-axis/diagonal couplings) then shifts and refines that mixing.

\paragraph{Decoupled limit ($g_{\mathrm{mix}}=0$).}
If the three axis channels are dynamically decoupled, each component supports an independent local phase choice
and an independent rank-1 Berry connection $a_\mu^{(x)},a_\mu^{(y)},a_\mu^{(z)}$.
In this limit, it is meaningful (and gauge-invariant) to speak of ``the Berry phase of an axis.''

\paragraph{Mixed regime ($g_{\mathrm{mix}}>0$).}
For any nonzero mixing, the physical eigenmodes are \emph{linear combinations} of axis labels.
Axis-specific phases become basis-dependent and are not gauge-invariant.
The correct adiabatic transport is instead a Wilczek--Zee connection on the transported near-degenerate subspace.
For an approximately threefold-degenerate carrier subspace, one obtains a $U(3)$ connection
$\mathcal{A}_\mu\in\mathfrak{u}(3)$ as in \Cref{eq:wz-connection}.
It can be decomposed into a trace $U(1)$ part and a traceless $SU(3)$ part,
\begin{equation}
  \mathcal{A}_\mu
  = \frac{1}{3}\,\mathrm{tr}(\mathcal{A}_\mu)\,I_3
  + \sum_{a=1}^{8} A^a_\mu\,T^a,
\end{equation}
where $T^a$ are $SU(3)$ generators.
This provides a concrete way to express how ``axis-aligned Berry phases'' transition into a non-Abelian color-like
connection once coupling is switched on.

\subsection{Berry curvature on a discretized Brillouin zone (link-variable method)}
\label{subsec:bz-berry}

For a lattice substrate, Berry curvature can be computed directly on the discretized Brillouin zone using
link variables and plaquettes, closely analogous to lattice-QCD Berry-curvature constructions
\citep{Fukui2005ChernDiscretizedBZ,PuYamamoto2018BerryLatticeQCD,Yamamoto2016BerryPhaseLatticeQCD}.
Let $\lvert u(\mathbf{k})\rangle$ be a normalized eigenvector (or $\mathbf{E}(\mathbf{k})$ a frame of a degenerate
subspace) on a momentum lattice with spacing $\tilde a$.
In the Abelian case one defines links
\begin{equation}
  U_j(\mathbf{k})=\frac{\langle u(\mathbf{k})\,|\,u(\mathbf{k}+\tilde{\jmath})\rangle}
  {|\langle u(\mathbf{k})\,|\,u(\mathbf{k}+\tilde{\jmath})\rangle|},
\end{equation}
and a plaquette
\begin{equation}
  U_{ij}(\mathbf{k}) = U_i(\mathbf{k})U_j(\mathbf{k}+\tilde{\imath})
  U_i^\dagger(\mathbf{k}+\tilde{\jmath})U_j^\dagger(\mathbf{k}),
\end{equation}
whose argument approximates the integrated Berry curvature over the plaquette.
In the non-Abelian case one uses overlap matrices between frames and projects to $U(n)$ links, then forms the same
plaquette structure to obtain a discretized non-Abelian curvature.
This is a natural diagnostic bridge between the lattice-substrate hypothesis and QCD-motivated non-Abelian geometry.
\subsection{Spinorial holonomy and the $4\pi$ aspect}

For $n=2$ (a two-level polarization subspace), holonomy lives naturally in a double cover of $SO(3)$.
A closed loop in physical space that induces a $2\pi$ rotation of an internal polarization frame can therefore
produce a sign change of the transported state, while a $4\pi$ loop returns it to itself.
In the present theory this spinorial $4\pi$ aspect is attributed to Berry/Wilczek--Zee holonomy of an internal
polarization bundle, and is treated as part of the particle-sector structure developed in \Cref{sec:localized}.
